\documentclass[11pt]{article}
\usepackage[margin=1in]{geometry}
\usepackage{amsmath,amssymb}
\usepackage{booktabs,array}
\usepackage{xurl}
\usepackage[hidelinks]{hyperref}
\setlength{\emergencystretch}{2em}
\title{Braking ends at rest: the target contract without reversal}
\author{}
\date{October 11, 2026}
\begin{document}
\maketitle
\paragraph{Status and scope.}
On October 11, 2026 the user decided that braking never turns a vehicle back:
``braking is braking; at most it brings the speed to zero''. The ego's
terminal law already ended at rest. The target contract did not: under the
signed-speed constant-$(A,\beta)$ forecast of
\path{predictiveSafetyGeometry.predictTarget}, a target whose speed reached
zero turned back along its track and accelerated in reverse. The simulated
truth followed the same forecast. This report records the change, its
checks, and a paired campaign against the braking-continuation terminal set
of commit \texttt{ef6ca13} (signed speed, otherwise identical).

\paragraph{Result.}
The forecast, the simulated truth and the terminal set's parked check now
stop a braking target at rest, and the full suite passes (608 of 608). The
campaign does not recover more encounters: 9 of 14 exact (9 before, the same
encounters) and 54 of 84 noisy (60 before). With a lead that comes to rest,
stopping behind it is now terminal. The 16-hold plan therefore follows a
braking lead in its lane instead of passing it. The declared model and the
plant cannot go below 1~m/s, so the plan fails as the lead slows to rest. In
the noisy runs, braking leads fall from 8 to 3 of 12. The change is the
physical contract the user asked for. It exposes the two limits named in the
previous report: the model cannot stop, and nothing in a short plan rewards
passing.

\section{The change}
\paragraph{Contract.}
The target keeps Sharma's constant tangential acceleration $A$ and sideslip
$\beta$. Its speed is $V+At$ until it reaches zero and zero afterwards. The
state's $A$ becomes 0 at rest, so the forecast keeps the semigroup property.
A target at rest moves off only forward ($A>0$). For a forward target the
speed is $\max(0,V+At)$ and the arc length is its integral
(\texttt{predictiveSafetyGeometry.targetArc}). A target braking from $V$
rests after $-V/A$ seconds at the arc length $-V^2/(2A)$. A target in reverse
($V<0$) that brakes ($A>0$) also stops.

\paragraph{Forecast epoch.}
The controller used to anchor the forecast at sample 0 by running the latest
estimate backward, $q_0=\phi(q,-kh)$, and forecast every later time from
$q_0$. A stop cannot be run backward: an accelerating target observed at
3~m/s after 10~s at $A=0.5$~m/s$^2$ has a backward speed of $-2$~m/s. Its
forward forecast from that epoch would then stop it after 4~s. The
controller now keeps the latest estimate and the sample index at which it
was taken (\texttt{targetEpoch}, \texttt{targetEpochIndex}; controller state
version 78) and forecasts forward from there. Running a forecast backward
through rest raises \texttt{collisionAvoidanceController:irreversibleTargetStop};
only the scenario construction runs backward, along motion that does not
pass through rest.

\paragraph{Terminal set.}
The parked check's swept stretch (\texttt{localTrack}) ends at the rest point
when the target comes to rest inside the range. Such a target never leaves,
and its rest pose must clear the parked ego. A target at rest without
forward acceleration sweeps a single point. The braking phase uses the
forecast directly.

\paragraph{Uncertainty enclosure.}
The observer's constant-parameter enclosure (\texttt{targetErrorTube},
\texttt{targetPositionSupport}) now uses the nominal arc of the contract.
Its arc-length error bound $\delta V\,|t|+\delta A\,t^2/2$ still holds. The
speed $\max(0,V+At)$ is 1-Lipschitz in $V+At$, so two forward targets'
arc lengths differ by at most the integral of $|\delta V+\delta A\,\tau|$.
A speed interval that straddles zero mixes forward and reverse targets, and
this bound does not cover that case.

\section{Checks}
\begin{itemize}
\item \emph{Unit tests.} The forecast tests now check the following.
  \begin{itemize}
  \item Braking ends at rest, at 1~m for 2~m/s and $-2$~m/s$^2$, and stays
    there, with the semigroup property and the batched form.
  \item A reverse-moving target that brakes stops.
  \item A target at rest with $A<0$ stays still; with $A>0$ it moves off forward.
  \item A turning target stops on its arc.
  \item Running backward through rest is refused, and a backward run that
    does not reach rest inverts exactly.
  \end{itemize}
  The terminal-set tests cover these cases.
  \begin{itemize}
  \item A braking target's swept stretch ends at its rest point and never
    leaves; at rest it is a point, or it moves off forward.
  \item A lead braking hard rests ahead of the parked ego with the clearance
    $10+4.5-6.4-4.8-0.2=3.1$~m.
  \item A lead that rests too close is met during the ego's braking.
  \end{itemize}
  Full suite: 608 of 608 (\path{tests.txt}; 606 before, plus two new tests).
\item \emph{Parked check against the forecast.}
  \path{scripts/checkParkedClearance.m} samples the target in time through
  \texttt{predictTarget} (200001 samples up to its first exit, its rest plus
  1~s, or one turn of its circle plus 1~s). This is independent of the
  closed form's swept stretch. The check covered 2895 cases (seed 11; 22
  needing more than 120~s skipped): 1216 braking, 646 sweeping a whole circle
  and 1826 never leaving the range. The closed form was never above the
  sampled least distance (at worst $5.2\cdot10^{-13}$~m, rounding) and at most
  $2.8\cdot10^{-7}$~m below it. The witness reproduced it to
  $5.2\cdot10^{-13}$~m, and the samples were at most 2.4~mm apart
  (\path{parked-clearance-check.json}).
\end{itemize}

\section{Campaign}
The 98 encounters of the earlier reports (seven geometries at 8 and 15~m/s;
seed 0 exact with the \texttt{ode45} plant, six seeds with the estimator in
the loop) were run for both arms: \texttt{ef6ca13} (signed speed) and the
working tree (stop at rest). Each arm ran three MATLAB sessions side by side
on the same encounter sequence. The signed arm reproduced the outcome and the
frame count of \path{report/BRAKING_TERMINAL_SET_20261011.tex} in all 98
encounters.

\begin{table}[h]
\centering\small
\caption{Paired campaign. No collision in either arm.}
\label{tab:campaign}
\begin{tabular}{lrrrr}
\toprule
 & \multicolumn{2}{c}{exact (14)} & \multicolumn{2}{c}{noisy (84)}\\
 & signed & stop & signed & stop\\
\midrule
recovered & 9 & 9 & 60 & 54\\
no solution: \texttt{nonlinearAcceptanceFailed} & 4 & 3 & 8 & 10\\
no solution: \texttt{anchorUnavailable} (plan ran out) & 1 & 2 & 2 & 2\\
no solution: \texttt{pcbfNoNumericalResult} & 0 & 0 & 12 & 16\\
no solution: \texttt{clfNoNumericalResult} & 0 & 0 & 2 & 2\\
least replay clearance (m) & 0.089 & 0.089 & 0.422 & 0.422\\
horizon, 95th percentile holds & 16 & 16 & 79 & 24\\
controller time per frame, 95th percentile (s) & 0.027 & 0.031 & 0.128 & 0.065\\
\bottomrule
\end{tabular}
\end{table}

\begin{table}[h]
\centering\small
\caption{Recovered encounters per geometry (both speeds).}
\label{tab:geometry}
\begin{tabular}{lrrrr}
\toprule
 & \multicolumn{2}{c}{exact (2 each)} & \multicolumn{2}{c}{noisy (12 each)}\\
 & signed & stop & signed & stop\\
\midrule
head-on & 2 & 2 & 11 & 11\\
accelerating head-on & 2 & 2 & 9 & 10\\
braking lead & 0 & 0 & 8 & 3\\
crossing & 1 & 1 & 12 & 10\\
turning crossing & 0 & 0 & 10 & 10\\
curved head-on & 2 & 2 & 6 & 7\\
curved crossing & 2 & 2 & 4 & 3\\
\bottomrule
\end{tabular}
\end{table}

The nine exact encounters that both arms recover are identical in frame
count, recovery time (median 9.9~s), lateral deviation and clearance. Before
its lead stops, a target's forecast and truth are the same under both
contracts. The unrecovered exact 8-m/s crossing and 15-m/s turning crossing,
whose targets do not brake, end 3 and 10 frames later than before. The
forecast is now evaluated from the latest estimate instead of through a
backward epoch, which changes it only by rounding, and these runs end at an
acceptance boundary. Of the 51 noisy encounters both arms recover, 42 have identical
frame counts. The noisy estimate's $A$ can be negative for a target that
does not brake, and its forecast then stops the target, so these runs
diverge. Three noisy encounters are newly recovered and nine newly lost
(\path{comparison.txt}, \path{campaign-runs.csv}, \path{failures.txt}).

\paragraph{Braking leads.}
The lead starts 6.25~m ahead at the ego's speed and brakes at 0.5~m/s$^2$.
It rests after 16~s (8~m/s) or 30~s (15~m/s). Under the signed contract it
then reverses toward the ego, and a parked ego behind a slow lead was not
terminal unless the lead's turning point lay outside the 50-m range. The
plan therefore had to pass. Noisy seeds 2 and 4 at 8~m/s passed it (largest
lateral position 5.6 and 5.2~m, least speed 5.6 and 6.2~m/s) and recovered
by 16.9 and 16.7~s. Under the stop contract, stopping behind the lead is
terminal. The same encounters followed the lead in its lane (largest lateral
position 1.3 and 1.4~m) and failed at 12.7 and 12.9~s at 1.32 and 1.69~m/s,
with the lead at 1.68 and 1.55~m/s. All six noisy 8-m/s braking leads end
this way, after 245--264 frames (\texttt{pcbfNoNumericalResult}), with an
endpoint stop time of 0.4--0.65~s. The exact braking leads fail later than
before, at 1.43~m/s after 13.2~s (8~m/s) and 2.72~m/s after 25.1~s (15~m/s),
against 3.58~m/s after 2.65~s and 1.28~m/s after 17.0~s.

\paragraph{Mechanism.}
The terminal set says the ego could brake to rest behind the lead. Nothing
in a 16-hold plan prefers passing to following, so the plan follows. The
declared model's domain ends at 1~m/s, as does the plant's, so the vehicle
cannot follow the lead to rest. Under the signed contract, the lead's
reversal had forced the pass that recovery requires.

\section{Decisions for the user}
\begin{enumerate}
\item \emph{A low-speed regime of the declared model and the plant.} Stopping
behind a stopped vehicle is now a terminal state that the vehicle cannot
reach. A kinematic bicycle blended in below a few m/s, with more integration
steps per hold, would let it.
\item \emph{A reason to pass.} A terminal cost on the endpoint (its CLF
value, or the progress along the path), or a minimum horizon longer than
the soft prefix, would let the planner prefer passing to following.
\item \emph{Keep the stop contract.} It is the user's decision and is
committed as the production contract; the signed contract is at
\texttt{ef6ca13}.
\end{enumerate}

\section{Reproduction}
\begin{itemize}
\item Unit tests: \texttt{matlab -batch "results = runtests('tests'); assertSuccess(results)"}.
\item Parked check: \path{scripts/checkParkedClearance.m} of this folder,
  \texttt{checkParkedClearance('<folder>/parked-clearance-check.json', 3000, 11)}.
\item Campaign: \path{scripts/run-workers.sh 2000 3} of this folder. The signed
  arm is a \texttt{git archive} of \texttt{ef6ca13}, the stop arm a copy of
  the working tree; each has a \texttt{solver} link. Then
  \texttt{python3 -I scripts/compare.py <out-signed> <out-stop> <folder>}.
  Raw outputs under \path{~/.cache/collisionAvoidance/target_stop_20261011/}
  are disposable and were deleted after the summaries were copied here.
\end{itemize}
\end{document}
